% !TeX program = lualatex
% Compile twice: lualatex 156.tex
% All references and prose are embedded; no BibTeX database or external inputs.
\documentclass[11pt,a4paper]{article}
\usepackage[margin=24mm,headheight=14pt]{geometry}
\usepackage{fontspec}
\setmainfont{Latin Modern Roman}
\setsansfont{Latin Modern Sans}
\setmonofont{Latin Modern Mono}
\usepackage{amsmath,amssymb,mathtools}
\usepackage{unicode-math}
\setmathfont{Latin Modern Math}
\usepackage{microtype}
\usepackage{enumitem,xurl,needspace}
\usepackage[dvipsnames]{xcolor}
\usepackage{hyperref}
\hypersetup{unicode=true,colorlinks=true,linkcolor=MidnightBlue,urlcolor=MidnightBlue,
 pdftitle={500 Mathematical Problem Statements},pdfauthor={Source-annotated compilation},bookmarksnumbered=true}
\usepackage{bookmark}
\usepackage{tocloft}
\setlength{\cftsecnumwidth}{3em}
\cftsetpnumwidth{2.5em}
\cftsetrmarg{4em}
\usepackage{titlesec}
\titleformat{\section}{\Large\bfseries\raggedright}{\thesection}{0.7em}{}
\usepackage{fancyhdr}
\pagestyle{fancy}\fancyhf{}
\fancyhead[L]{\small\sffamily Mathematical problem statements}
\fancyhead[R]{\small Source edition 22}
\fancyfoot[C]{\thepage}
\setlength{\parindent}{0pt}
\setlength{\parskip}{5pt plus 1pt}
\setlength{\emergencystretch}{3em}
\setcounter{tocdepth}{1}
\setcounter{secnumdepth}{1}
\setlist[itemize]{leftmargin=3em,itemsep=5pt,topsep=4pt,parsep=0pt}
\allowdisplaybreaks
\newcommand{\N}{\mathbb{N}}
\newcommand{\Z}{\mathbb{Z}}
\newcommand{\Q}{\mathbb{Q}}
\newcommand{\R}{\mathbb{R}}
\newcommand{\C}{\mathbb{C}}
\newcommand{\F}{\mathbb{F}}
\newcommand{\A}{\mathbb{A}}
\newcommand{\PP}{\mathbb P}
\newcommand{\E}{\mathbb{E}}
\newcommand{\eps}{\varepsilon}
\newcommand{\dd}{\,\mathrm d}
\newcommand{\sref}[2]{\hyperlink{source:#1:#2}{[S#2]}}
\newcommand{\eref}[2]{\hyperlink{extra:#1:#2}{[E#2]}}
% Turn the next switch off for a shorter reading copy. The quotations preserve
% every exactTarget field, including qualifications omitted from a short summary.
\newif\ifshowcatalogue
\showcataloguetrue
\newcommand{\cataloguescope}[1]{\ifshowcatalogue
 \paragraph{Exact scope in the supplied catalogue.}
 \begin{quote}\small #1\end{quote}\fi}

% Standalone research-notebook wrapper; the source preamble above is retained.
\numberwithin{equation}{section}
\hypersetup{pdftitle={Research attempt -- problem 156},
 pdfauthor={Research notebook; original compilation retained}}
\fancyhead[L]{\small\sffamily Research notebook}
\fancyhead[R]{\small\sffamily Problem 156 | 27 September 2026}
\begin{document}
\setcounter{section}{155}
% ================================================================
% problemId: problem.leopoldt-s-conjecture
% source releaseRank: 156; source releaseStatus: open_with_solved_subcases
% exactTarget SHA-256: 432b21df2eb77d5b50a40e702978d386173e03a1e84b8f5ddeee070303b49e6e
\clearpage
\section{\texorpdfstring{Leopoldt's conjecture}{Leopoldt's conjecture}}\label{156}
\subsection{Definitions and mathematical statement}
Let $\mathcal O_K^\times$ be the units of a number field $K$, of rank $r_1+r_2-1$. For a prime $p$, use the local $p$-adic logarithms to define
\[
 \ell_p:\mathcal O_K^\times\otimes_{\Z}\Q_p\longrightarrow\prod_{v\mid p}K_v,
 \qquad u\longmapsto(\log_v u)_v.
\]
The logarithm on units is extended from principal units and kills torsion. The Leopoldt defect is
$\delta_{K,p}=r_1+r_2-1-\dim_{\Q_p}\operatorname{im}\ell_p$.
Leopoldt's conjecture asserts $\delta_{K,p}=0$ for every $K,p$, equivalently injectivity of $\ell_p$. This is the full-rank assertion behind the $p$-adic regulator, not a statement that the ordinary real regulator vanishes or that local logarithms are algebraically independent.
\par\smallskip\textit{Formulation and concept references:} \sref{156}{1}.
\cataloguescope{Prove or disprove that for every number field K and every prime p, the Leopoldt defect of K at p is zero, equivalently that the p-adic regulator has the conjectured full rank.}
\subsection{Short English statement}
Do independent global units remain linearly independent when measured by their logarithms in all completions above any chosen prime?
% BEGIN RESEARCH ADDITION ATTEMPT-156
\subsection{Research attempt}

\paragraph{Current frontier.}
Ferri and Johnston's 2026 paper (version details in \eref{156}{1}) proves representation-theoretic restrictions
on defects and constructs infinite nonabelian families satisfying Leopoldt
at prescribed finite sets of primes. Of particular use below is their
Corollary 1.11: for an $S_3$-extension $L/\Q$ and a cubic subfield $K$,
$\delta_{L,p}=2\delta_{K,p}$ \sref{156}{1}. These results do not establish
injectivity for every number field and every prime. The attempt below gives
an explicit congruence family and a complete finite-residue calculation at
$p=3$. It is a special-case lemma, not a claimed improvement over all known
families.

\paragraph{Attack route.}
Representation descent can transfer a regulator result from a cubic field
to its normal closure, but does not eliminate the cubic input. A
transcendence route must prove independence over $\Q_p$, not merely rule out
rational relations among logarithms. We instead look for two explicit units
whose suitably scaled logarithms have independent reductions modulo $3$.
A nonzero residue minor proves exact $3$-adic independence: numerical
precision is then a certificate rather than a heuristic. The congruence
transfer idea is related to the explicit-unit methods reviewed in
\sref{156}{1}, Theorem 6.1 and the section on infinite families. All
calculations for the following polynomial are supplied here.

\paragraph{Attempt: an explicit family.}
\textbf{Lemma candidate.} For every integer $a\ge5$ with $a\equiv5\pmod9$,
let $\theta$ be a root of
\begin{equation}\label{eq156-family}
 f_a(T)=T^3-aT^2+(a-1)T+1,
\end{equation}
let $K_a=\Q(\theta)$, and let $L_a$ be its Galois closure. Then $K_a$ is
totally real of degree three, $\operatorname{Gal}(L_a/\Q)\cong S_3$, and
\[
 \delta_{K_a,3}=\delta_{L_a,3}=0.
\]
Moreover, this parameter set contains infinitely many isomorphism classes
of fields $K_a$, and hence infinitely many normal closures $L_a$.

\textit{Algebraic preliminaries.} A rational root of the monic cubic must be
$1$ or $-1$. But
\[
 f_a(1)=1,\quad f_a(-1)=1-2a,\quad f_a(0)=1,\quad
 f_a(2)=7-2a,\quad f_a(a)=a^2-a+1.
\]
There are therefore three distinct real roots, one in each of
$(-1,0)$, $(1,2)$, and $(2,a)$, and no rational root. In particular the
unit rank is two. The polynomial discriminant is
\begin{equation}\label{eq156-discriminant}
 D(a)=a^4-2a^3-5a^2+6a-23=(a^2-a-3)^2-32.
\end{equation}
Write $B=a^2-a-3$. For $a\ge5$ we have $B\ge17$ and
\[
 (B-1)^2<D(a)<B^2,
\]
since $D(a)-(B-1)^2=2B-33>0$. Hence $D(a)$ is not a square. The
irreducible cubic has Galois group $S_3$. Both
\[
 u=\theta,\qquad v=\theta-1
\]
are units, because $N_{K_a/\Q}(u)=N_{K_a/\Q}(v)=-1$.

\textit{The local algebra.} For $a\equiv5\pmod9$, reduction modulo three
gives
\[
 f_a(T)\equiv(T+1)(T^2+1),\qquad D(a)\equiv2\pmod3.
\]
Thus $3$ does not divide the index $[\mathcal O_{K_a}:\Z[\theta]]$, and
\[
 A_a:=\Z_3[T]/(f_a)\cong\mathcal O_{K_a}\otimes\Z_3
          \cong\prod_{w\mid3}\mathcal O_{K_a,w}.
\]
In particular $1,\theta,\theta^2$ are an integral $\Z_3$-basis of this
product algebra. The norm used here is the $3$-adic one, not a single
archimedean embedding.

\textit{A residue calculation with two units.} In the quotient by $f_5$,
ordinary polynomial division yields the exact identities
\[
 T^8\equiv5082T^2-5523T-1298,\qquad
 (T-1)^8\equiv480T^2-525T-119.
\]
Since $f_a\equiv f_5\pmod9$, these imply in $A_a/9A_a$
\[
 u^8\equiv7+3\theta+6\theta^2,
 \qquad v^8\equiv7+6\theta+3\theta^2.
\]
Both eighth powers belong to $1+3A_a$. Define
$q_u=(u^8-1)/3$ and $q_v=(v^8-1)/3$ in $A_a$. Then
\begin{equation}\label{eq156-residues}
 q_u\equiv2+\theta+2\theta^2,\qquad
 q_v\equiv2+2\theta+\theta^2\pmod{3A_a}.
\end{equation}
This is a statement for the whole congruence family, not a computation at
only one value of $a$.

\textit{From residues to exact logarithms.} For every $q\in A_a$,
\[
 \frac13\log(1+3q)
 =q+\sum_{k\ge2}\frac{(-1)^{k+1}3^{k-1}q^k}{k}
 \equiv q\pmod{3A_a}.
\]
The congruence and convergence are justified term by term: for $k\ge2$,
$k-1-v_3(k)\ge1$, and these valuations tend to infinity. The local
logarithms in the question satisfy $\log(u^8)=8\log u$, and similarly for
$v$. Consequently the two integral vectors
\[
 \frac83\ell_3(u),\qquad\frac83\ell_3(v)
\]
have, in the displayed integral basis, the reduction matrix
\begin{equation}\label{eq156-minor}
 \begin{pmatrix}2&2\\1&2\\2&1\end{pmatrix}\quad\text{over }\F_3.
\end{equation}
Its first two rows have determinant $2\ne0$ in $\F_3$. The corresponding
minor over $\Z_3$ is therefore a unit. The two logarithm vectors are
linearly independent over $\Q_3$.

It follows at once that $u$ and $v$ are multiplicatively independent modulo
torsion. Since the global unit rank is two, they span
$\mathcal O_{K_a}^{\times}\otimes\Q_3$: it is unnecessary to show that they
are a fundamental integral system of units. Thus $\ell_3$ is injective on
that entire vector space, proving $\delta_{K_a,3}=0$. Finally apply the
$S_3$ defect identity from \sref{156}{1}, Corollary 1.11, to obtain
$\delta_{L_a,3}=2\delta_{K_a,3}=0$. The transfer is an established theorem;
the residue-minor computation above is the deduction made in this attempt.

\paragraph{Attempt: infinitude of distinct fields, not just parameters.}
We include this argument because an infinite polynomial family can in
principle repeat finitely many number fields. First, any nonconstant
integer polynomial $Q$ has infinitely many prime divisors of its nonzero
integer values. Here is the elementary proof needed. If those primes were
the finite set $\mathcal P$, choose $n_0$ with $Q(n_0)\ne0$, put
$P=\prod_{\ell\in\mathcal P}\ell$, and consider
\[
 \frac{Q(n_0+Q(n_0)Pt)}{Q(n_0)}.
\]
This is an integer polynomial in $t$, congruent to $1$ modulo every
$\ell\in\mathcal P$, and unbounded in absolute value. Its values could
have no prime divisors, although some have absolute value greater than one,
a contradiction.

Apply this fact to $Q(n)=D(9n+5)$. The polynomial $D$ has no repeated roots
in characteristic zero: from
$D'=2(a^2-a-3)(2a-1)$ a common root would either have
$D=-32$, or be $a=1/2$, where $D=-343/16$. Thus, outside finitely many
primes, a root of $D$ modulo $\ell$ is simple. Infinitely many primes
$\ell\ne3$ therefore divide some $D(a_0)$ with $a_0\equiv5\pmod9$ and
$D'(a_0)\not\equiv0\pmod\ell$.

For such a prime replace $a_0$ by $a=a_0+9\ell t$. Taylor expansion in
integers gives
\[
 D(a)\equiv D(a_0)+9\ell tD'(a_0)\pmod{\ell^2}.
\]
Exactly one residue class of $t$ modulo $\ell$ makes this vanish modulo
$\ell^2$. Choose another class and then a sufficiently large representative
so that $a\ge5$. It follows that $v_\ell(D(a))=1$. Since
\[
 D(a)=\operatorname{disc}(K_a)
       [\mathcal O_{K_a}:\Z[\theta]]^2,
\]
the prime $\ell$ ramifies in $K_a$. A finite list of number fields has only
finitely many ramified rational primes. The construction supplies
infinitely many possible ramified primes, and hence infinitely many
isomorphism classes among the $K_a$. A fixed degree-six normal closure has
only finitely many subfields, so the same conclusion holds for the $L_a$.
This completes the proof of the stated family lemma. \hfill$\square$

\paragraph{Stress test.}
The prime $3$ is unramified here; otherwise using $1,\theta,\theta^2$ as a
basis of the full local integer algebra would require a new index check.
The factor $8/3$ is legitimate because the eighth powers are in $1+3A_a$;
we have not applied a principal-unit logarithm series directly to $\theta$.
The nonzero minor is an exact residue obstruction to dependence, not a
floating-point determinant that might be numerically small. Conversely,
a zero minor at this precision would not prove a positive defect.

At $a=4$ the discriminant is $49$, so the $S_3$ assertion would fail;
that parameter is excluded. Nothing in the argument handles an arbitrary
prime $p\ne3$, or a general number field without these explicit units.
The proposed family is not asserted to be novel in the literature, and the
existence of infinite nonabelian families for finite prime sets is already
part of \sref{156}{1}. What is recorded as a lemma candidate is the
specific polynomial, congruence, and certificate proved above. The checking
script repeats the polynomial divisions and residue tests exactly; the
symbolic proof, not an enumeration of parameters, supplies the universal
quantifier over $a\equiv5\pmod9$.

\paragraph{Outcome.}
\textbf{Outcome: New partial result / lemma candidate.}

The attempt obtains an explicit, fully argued family statement:
Leopoldt holds at $3$ for the totally real cubic fields
\eqref{eq156-family} with $a\ge5$, $a\equiv5\pmod9$, and for their
$S_3$ normal closures. There are infinitely many distinct such fields.
The label records a lemma derived in this notebook, not a verified claim
of literature priority.

A complete resolution would require a mechanism forcing full logarithmic
rank for every field and every prime, rather than this particular nonzero
residue minor. Representation-theoretic transfer does not create the
missing independent logarithms. No general defect bound is improved and no
counterexample to Leopoldt is produced.
% END RESEARCH ADDITION ATTEMPT-156
\subsection{Sources}
\begin{itemize}\raggedright
\item[\textnormal{[S1]}]\hypertarget{source:156:1}{} Fabio Ferri and Henri Johnston, Research in Number Theory 12 (2026), article 32.
\url{https://link.springer.com/article/10.1007/s40993-026-00717-2}.
{\small\textit{Catalogue formulation source.}}
% BEGIN RESEARCH ADDITION SOURCES-156
\item[\textnormal{[E1]}]\hypertarget{extra:156:1}{} Fabio Ferri and Henri
Johnston, \emph{Applications of representation theory and of explicit units
to Leopoldt's conjecture}, arXiv:2301.05700v2, 23 March 2026;
\emph{Research in Number Theory} 12 (2026), article 32.
\url{https://arxiv.org/abs/2301.05700v2}.
{\small\textit{Full-title/version identification of the paper retained as
[S1], not an independent additional result. Corollary 1.11 supplies the
$S_3$ transfer; Theorem 6.1 explains the finite-residue strategy.
Consulted 27 September 2026.}}
% END RESEARCH ADDITION SOURCES-156
\end{itemize}

\end{document}
